\subsection{Over-regularized warmup}

Set
\begin{equation}
  \bar\rho:=(1+\sqrt\alpha)\rho.
  \label{eq:overrho}
\end{equation}
The proof-oriented hybrid template is:
\begin{enumerate}[label=\arabic*.,leftmargin=*]
  \item run an accelerated local method on $F_{\bar\rho}$ until a certified
  coarse objective gap is reached;
  \item retract the accelerated point to a lower certificate
  $\ell\le x_{\bar\rho}^\star$ and discard the momentum state; and
  \item run a one-sided $\omega=1$ local Gauss--Seidel/proximal coordinate
  tail on $F_\rho$.
\end{enumerate}
The separation between the accelerated point and the retracted checkpoint is
essential.  Momentum is not itself a valid warm-start invariant for a
monotone residual process.

The choice~\eqref{eq:overrho} creates KKT slack of scale
$\alpha(\bar\rho-\rho)=\alpha^{3/2}\rho$.  This is large enough to make the
tail cheap, but changes the support budget only by a $1+O(\sqrt\alpha)$
factor.

\subsection{Why the tail uses \texorpdfstring{$\omega=1$}{omega=1}}

For a quadratic coordinate update, the exact objective change is negative for
every fixed relaxation $0<\omega<2$.  Objective descent alone is not the
locality invariant needed here: when $\omega>1$, signed residuals may be
created and must be charged separately.  The signed warm-start SOR lemma used
in the broader project allows $\omega\in[1,2]$ under its own hypotheses, but
the unconditional one-sided mass argument below uses the safe endpoint
$\omega=1$.  The note therefore makes no claim that the optimized SOR
parameter preserves the same tail bound.

\subsection{The desired division of work}

Let $T=\softO(1/\sqrt\alpha)$ denote the number of accelerated proximal
centers needed for the warmup gap.  The intended theorem decomposes as
\[
  W_{\mathrm{hybrid}}
  =W_{\mathrm{warm}}+W_{\mathrm{tail}}.
\]
The tail term is already controlled at the target scale.  Every unsuccessful
proof in this note fails at the same place: bounding $W_{\mathrm{warm}}$
without silently charging the whole graph or assuming monotone accelerated
support.
